\section{Cluster y su verificación}
\label{sec:cluster-verificacion}

\subsection{Levantamiento del clúster de tres nodos}

El clúster se define en \texttt{db-cluster/docker-compose.cockroach.yml} con tres servicios
(\texttt{roach1}, \texttt{roach2}, \texttt{roach3}) sobre la red \texttt{bridge}
\texttt{soporte-net}, cada uno con volumen persistente propio y el comando
\texttt{start --insecure --join=roach1,roach2,roach3} apuntando al resto de nodos —siguiendo el
patrón estándar de arranque de un clúster CockroachDB multi-nodo. Cada nodo declara además una
\texttt{locality} distinta (\texttt{quevedo-centro}, \texttt{quevedo-norte},
\texttt{quevedo-sur}), simulando las tres zonas operativas del ISP para fines de política de
replicación, y expone un \emph{healthcheck} HTTP sobre \texttt{/health?ready=1} que Docker Compose
usa para reportar el estado del contenedor.

\begin{lstlisting}[language=bash, caption={Extracto de \texttt{docker-compose.cockroach.yml}: definición de un nodo}, label={lst:compose-crdb}]
roach1:
  image: cockroachdb/cockroach:latest-v23.2
  command: start --insecure --max-offset=4500ms
    --locality=zone=quevedo-centro
    --join=roach1,roach2,roach3
  ports: ["26257:26257", "8083:8080"]
  volumes: ["roach1-data:/cockroach/cockroach-data"]
  networks: [soporte-net]
\end{lstlisting}

Tras \texttt{docker compose -f docker-compose.cockroach.yml up -d}, el clúster se inicializa una
única vez con \texttt{cockroach init --insecure} ejecutado contra cualquiera de los tres nodos.
La verificación de que los tres nodos forman efectivamente un único clúster operativo se realiza
con \texttt{cockroach node status --insecure}, que debe reportar los tres identificadores de nodo
con \texttt{is\_live=true} e \texttt{is\_available=true}; el mismo estado es visible en la consola
web del clúster (puerto 8083 en este despliegue, reasignado desde el 8080 por defecto para no
chocar con otro servicio local). Sobre este clúster ya verificado operativo se aplica el esquema
fragmentado descrito en la Sección~\ref{sec:diseno-esquema} y se ejecutan las pruebas de
integración y la instrumentación de métricas de las siguientes subsecciones.

\subsection{Pruebas de integración contra CockroachDB real}

Para evitar que la corrección del acceso a datos dependiera únicamente de pruebas unitarias con
\emph{mocks} —que no pueden reproducir comportamiento específico de CockroachDB como el
aislamiento \texttt{SERIALIZABLE} o los tipos nativos \texttt{UUID}/\texttt{STRING}—, se añadió una
suite de pruebas de integración basada en Testcontainers
(\texttt{TicketRepositoryIntegrationTest.java}) que levanta un contenedor real de
\texttt{cockroachdb/cockroach:latest-v23.2} en modo \texttt{start-single-node} para cada ejecución
de la suite y lo descarta automáticamente al finalizar (limpieza gestionada por el
\emph{sidecar} Ryuk de Testcontainers, sin intervención manual). El esquema ya no lo crea la
prueba a mano: con \texttt{flyway-core} en el \emph{classpath}, Spring Boot aplica
automáticamente la migración real
\texttt{V1\_\_init\_ticket\_schema.sql} (Entregable 5 de la guía de cierre) contra esa misma base
al construir el contexto de la aplicación —el guion suelto \texttt{db-cluster/scripts/init\_db.sql}
que esta prueba citaba antes de esa corrección ya no existe en el repositorio, precisamente para
que una migración rota se detecte aquí en vez de poder divergir en silencio de un guion aparte.

La suite contiene dos pruebas. La primera ejerce el mismo \texttt{TicketRepository} usado en
producción con una operación real de guardado y búsqueda por \texttt{findByTicketId(UUID)} —el
mismo camino que ejercita el índice único secundario \texttt{tickets\_id\_key} descrito en la
Sección~\ref{sec:diseno-esquema}—. La segunda es una prueba empírica, no simulada, de que el
clúster aplica de verdad aislamiento \texttt{SERIALIZABLE}: dos hilos abren cada uno su propia
conexión y transacción, ambos leen la misma fila de \texttt{tickets} (sincronizados con una
barrera para garantizar que ambas lecturas ocurren antes de cualquier escritura) y luego ambos
intentan actualizar el mismo campo \texttt{status}. El resultado esperado —y observado— es que al
menos una de las dos transacciones falle al hacer \texttt{commit} con un error de conflicto de
escritura serializable, identificado por el código de estado SQL \texttt{40001}:

\begin{lstlisting}[language=Java, caption={Extracto de la prueba de conflicto serializable}, label={lst:testcontainers}]
assertThat((Exception) conflict.get())
    .describedAs("al menos una de las dos transacciones concurrentes "
        + "debio fallar por conflicto serializable")
    .isNotNull();
assertThat(conflict.get().getSQLState()).isEqualTo("40001");
\end{lstlisting}

Esta prueba corrobora en código, de forma reproducible en cualquier máquina con Docker
(incluyendo el entorno de integración continua), el mismo tipo de error observado en producción
en \texttt{report-service} (\texttt{ReportEventListener.applyWithRetry}) y que
\texttt{TicketService} está preparado para reintentar automáticamente (ver
Sección~\ref{sec:cluster-verificacion}, métricas de reintento). Ambas pruebas se ejecutaron de
forma exitosa contra la imagen real de CockroachDB antes de la integración de esta sección.

\subsection{Instrumentación de métricas operativas}

\texttt{ticket-service} expone tres métricas Prometheus, registradas mediante Micrometer y
consumibles en el \emph{endpoint} \texttt{/actuator/prometheus}, descritas en la
Tabla~\ref{tab:metricas}.

\begin{table}[H]
\centering
\caption{Métricas Prometheus expuestas por \texttt{ticket-service}}
\label{tab:metricas}
\begin{tabular}{@{}p{4.3cm}p{1.8cm}p{6.7cm}@{}}
\toprule
\textbf{Métrica} & \textbf{Tipo} & \textbf{Significado} \\
\midrule
\texttt{crdb\_query\_duration\_seconds} & Histograma & Duración de cada consulta a \texttt{TicketRepository}, con histograma de percentiles \\
\texttt{crdb\_transaction\_retries\_total} & Contador & Reintentos por conflicto de escritura serializable \\
\texttt{crdb\_pool\_active\_connections} & \emph{Gauge} & Conexiones activas del \emph{pool} HikariCP hacia CockroachDB \\
\bottomrule
\end{tabular}
\end{table}

El histograma de duración se alimenta desde un \emph{aspect} de Spring AOP
(\texttt{RepositoryTimingAspect}) que envuelve cada llamada al repositorio sin requerir
instrumentación manual en cada método de servicio. El contador de reintentos se incrementa
explícitamente desde \texttt{TicketService} cada vez que una operación de guardado captura una
\texttt{ConcurrencyFailureException} y reintenta —el mismo mecanismo ejercitado por la prueba de
conflicto serializable de la subsección anterior—. El \emph{gauge} de conexiones activas no
mantiene un valor cacheado: Micrometer lo reconsulta bajo demanda directamente contra el
\texttt{HikariPoolMXBean} real del \emph{pool}, por lo que siempre refleja el estado actual del
\emph{pool}, no una fotografía tomada en el momento del registro.

Las tres métricas se validaron con el \emph{linter} oficial de Prometheus (\texttt{promtool check
metrics}) ejecutado sobre la salida cruda del \emph{endpoint}, dentro de un contenedor
\texttt{prom/prometheus:latest} para no requerir una instalación local de Prometheus. La
validación no reportó advertencias ni errores sobre los tres instrumentos: los nombres siguen la
convención de sufijo por unidad (\texttt{\_seconds}, \texttt{\_total}), y las descripciones
(\texttt{HELP}) y tipos (\texttt{TYPE}) están correctamente declarados en la exposición.

\subsection{Tolerancia a fallos del clúster}

El protocolo de medición de tolerancia a fallos consiste en generar una carga de escritura
sostenida contra el clúster de tres nodos, derribar deliberadamente un nodo
(\texttt{docker stop roach2}) y registrar la latencia P50/P95 de las operaciones de escritura
antes, durante y después de la caída, repitiendo después el mismo procedimiento derribando dos
nodos simultáneamente para observar la pérdida de disponibilidad de escritura predicha por la
teoría de consenso de la Sección~\ref{sec:fundamento-teorico} (con dos de tres nodos caídos ya no
existe mayoría para confirmar escrituras vía Raft).

Durante la preparación de este protocolo se recreó, de forma no planificada pero informativa, el
escenario adyacente de saturación de memoria: una carga inicial de 150{,}000 filas en una sola
transacción provocó que el nodo \texttt{roach1} fuera terminado por el sistema operativo por falta
de memoria (código de salida 137) bajo el límite original de 512\,MB por contenedor, y el propio
reinicio del nodo —una operación intensiva en memoria en CockroachDB, que debe reconstruir su
estado desde el registro de Raft— volvió a fallar por el mismo motivo antes de estabilizarse. El
límite de memoria de los tres nodos se incrementó a 1024\,MB en
\texttt{docker-compose.cockroach.yml} como mitigación documentada, y la recarga completa del
conjunto de datos de 150{,}000 filas se replanificó en lotes más pequeños para el protocolo formal
de tolerancia a fallos. Este incidente, aunque no forma parte del protocolo experimental
diseñado, es evidencia operativa consistente con la teoría: un nodo bajo presión de recursos deja
de poder participar en el consenso, degradando la disponibilidad del clúster de la misma forma
que una caída deliberada, y su recuperación no es instantánea.

\subsubsection{Resultados}

El protocolo se ejecutó con una carga inicial de 101{,}996 filas en \texttt{tickets}. La
Tabla~\ref{tab:tolerancia-1nodo} resume la carga sostenida (\texttt{load\_write.py},
100 filas/s) durante la caída de \texttt{roach2}.

\begin{table}[H]
\centering
\caption{Latencia por ventana de 10s durante la caída de 1 nodo (roach2)}
\label{tab:tolerancia-1nodo}
\begin{tabular}{@{}lrrrr@{}}
\toprule
\textbf{Ventana} & \textbf{P50} & \textbf{P95} & \textbf{Errores} & \textbf{Filas} \\
\midrule
1 & 38.6 ms & 118.4 ms & 0 & 110 \\
2 (\texttt{docker kill roach2}) & 69.3 ms & \textbf{1590.8 ms} & 0 & 44 \\
3 & 31.0 ms & 107.9 ms & 0 & 130 \\
4 & 40.2 ms & 95.7 ms & 0 & 114 \\
5 & 26.3 ms & 69.3 ms & 0 & 163 \\
\bottomrule
\end{tabular}
\end{table}

El pico de $P95=1590.8$\,ms en la ventana 2 coincide con el instante de \texttt{docker kill
roach2}: el tiempo que Raft necesita para reelegir líder en los rangos cuyo \emph{leaseholder} era
roach2. Ningún dato se pierde ni falla (0 errores en las 691 filas insertadas); el único efecto
observable es latencia transitoria, no indisponibilidad — el clúster tolera la caída de un nodo
sin degradar la disponibilidad de escritura, como predice el quórum de mayoría 2 de 3.

Al repetir la prueba derribando también \texttt{roach3} (solo \texttt{roach1} en pie), la carga
sostenida (50 filas/s) muestra una brecha de \textbf{62 segundos} entre dos ventanas consecutivas
que normalmente distan 10-12\,s, coincidiendo exactamente con la ventana de tiempo en la que
ambos nodos permanecieron caídos: sin mayoría posible, la conexión abierta quedó bloqueada
esperando confirmación en vez de recibir un error inmediato, y una consulta
\texttt{SELECT count(*)} lanzada en ese mismo intervalo no devolvió resultado hasta reincorporar
\texttt{roach2}. Esta indisponibilidad por bloqueo —en vez de una excepción explícita del
cliente— es una manifestación válida de la pérdida de disponibilidad de escritura predicha por la
teoría de consenso: con un solo nodo vivo de tres no existe mayoría para que Raft confirme ninguna
escritura, exactamente el límite documentado en la Sección~\ref{sec:fundamento-teorico}. La
recuperación es inmediata al ejecutar \texttt{docker start roach2}: con dos nodos vivos vuelve a
haber quórum y la latencia regresa a valores normales (20.6/43.9\,ms). El detalle completo,
incluyendo la secuencia exacta de comandos, está en la bitácora
(\texttt{docs/evidencias/tolerancia\_fallos.md}) y el video
(\texttt{docs/evidencias/tolerancia\_fallos.mp4}, 4:55~min).

\begin{figure}[H]
\centering
\includegraphics[width=0.8\textwidth]{figuras/consola-recuperacion.jpg}
\caption{Consola web de CockroachDB tras la recuperación final: los tres nodos en estado
\texttt{LIVE}, 61 réplicas cada uno (captura del video de tolerancia a fallos).}
\label{fig:consola-recuperacion}
\end{figure}
